\documentclass[11pt,a4paper]{article}

\usepackage[utf8]{inputenc}
\usepackage[T1]{fontenc}
\usepackage{lmodern}
\usepackage{geometry}
\usepackage{amsmath,amssymb,amsfonts,bm}
\usepackage{booktabs}
\usepackage{array}
\usepackage{enumitem}
\usepackage{graphicx}
\usepackage{xcolor}
\usepackage{hyperref}
\usepackage[numbers,sort&compress]{natbib}

\geometry{margin=2.5cm}
\setlist[itemize]{topsep=2pt,itemsep=2pt}
\setlist[enumerate]{topsep=2pt,itemsep=2pt}

\hypersetup{
    colorlinks=true,
    linkcolor=blue!50!black,
    citecolor=blue!50!black,
    urlcolor=blue!50!black
}

\newcommand{\R}{\mathbb{R}}
\newcommand{\N}{\mathbb{N}}
\newcommand{\calD}{\mathcal{D}}
\newcommand{\calE}{\mathcal{E}}
\newcommand{\calG}{\mathcal{G}}
\newcommand{\calK}{\mathcal{K}}
\newcommand{\calN}{\mathcal{N}}
\newcommand{\calP}{\mathcal{P}}
\newcommand{\calR}{\mathcal{R}}
\newcommand{\calS}{\mathcal{S}}
\newcommand{\calT}{\mathcal{T}}
\newcommand{\calV}{\mathcal{V}}
\newcommand{\KNN}{\operatorname{kNN}}
\newcommand{\rdeg}{\operatorname{rdeg}}
\newcommand{\argmax}{\operatorname*{arg\,max}}
\newcommand{\argmin}{\operatorname*{arg\,min}}

\graphicspath{{figures/}{figures/2018/}{figures/generated/}}

\newcommand{\placeholderfigure}[2]{%
\begin{center}
\fbox{%
\begin{minipage}[c][0.22\textheight][c]{0.86\linewidth}
\centering
\vspace{0.3em}
\textbf{Figure placeholder}\par
\vspace{0.5em}
#1\par
\vspace{0.5em}
\small Expected file: \texttt{#2}
\end{minipage}}
\end{center}}

\newcommand{\maybegraphics}[3][]{%
\IfFileExists{#2}{\includegraphics[#1]{#2}}{\placeholderfigure{#3}{#2}}}

\title{Round-Trip Closure on Directed kNN Graphs for Multiscale Clustering}
\author{Eraldo Pereira Marinho\\
\small Department of Statistics, Applied Mathematics and Computing\\
\small S\~ao Paulo State University (UNESP), Rio Claro, Brazil}
\date{}

\begin{document}

\maketitle

\begin{abstract}
We introduce a graph-based clustering formulation based on recursive round-trip closure in directed $k$-nearest-neighbor graphs. The method originates from an earlier relational model in which an asymmetric relation $T\subseteq A\times B$ induces a reflective operator $R_T=T^{-1}\circ T$. In the present formulation, this relational mechanism is instantiated as a directed kNN graph over a finite dataset. Local clusters are obtained by iterating a round-trip neighborhood operator until saturation, while the evolution of the resulting partitions across increasing values of $K$ provides a multiscale description of cluster formation, persistence, and merging. The method combines directed kNN topology, reverse-degree popularity, optional visibility constraints, saturation depth, and $K$-layer transition analysis. It should be understood as a new variant within the broader family of nearest-neighbor and graph-based clustering methods, rather than as a generic first use of kNN graphs for clustering.
\end{abstract}

\noindent\textbf{Keywords:} graph-based clustering; k-nearest-neighbor graph; round-trip closure; shared nearest neighbors; mutual kNN; reverse nearest neighbors; multiscale clustering; saturation.

\section{Introduction}

Nearest-neighbor relations are among the most natural structures for representing finite datasets embedded in metric spaces. They appear in classification, nonparametric density estimation, graph construction, dimensionality reduction, outlier detection, and clustering \citep{CoverHart1967,LoftsgaardenQuesenberry1965,Bentley1975,FriedmanBentleyFinkel1977,VonLuxburg2007}. In clustering, nearest-neighbor information has been used through shared neighbors, reciprocal neighbors, mutual kNN graphs, reverse nearest neighbors, and density-based criteria \citep{JarvisPatrick1973,ErtozSteinbachKumar2003,BritoChavezQuirozYukich1997,GowdaKrishna1978,ChowdhuryAmorim2018,XieLeeWangChenZhou2019}.

This paper reconstructs and extends a line of work initiated in an unpublished 2013 relational note on redundancy-based information transport \citep{Marinho2013RedundanceRevised} and further developed in an unpublished kNN redundancy-network manuscript \citep{Marinho2018KNN}. The key idea is that a directed relation can be used not only to connect points, but also to define an operator of reflective expansion. A point first reaches its directly related objects; then the inverse relation brings back all objects sharing the same targets. Repeating this procedure yields a closure process that grows until saturation.

In the current setting, the abstract relation is instantiated as a directed kNN graph. Given a dataset $X=\{x_1,\ldots,x_N\}\subset\R^d$, each point $x_i$ points to its $K$ nearest neighbors. This directed graph induces both a forward neighborhood and a reverse neighborhood. The proposed clustering mechanism uses this directed topology to perform a recursive round-trip expansion. The expansion stops when no additional points are reached. The saturated sets, their overlaps, and their organization across $K$ define the cluster structure.

The contribution is not the isolated use of kNN, shared neighbors, mutuality, or reverse degree. These ingredients are well established. The contribution is the combination of:

\begin{enumerate}
    \item a recursive round-trip closure operator on a directed kNN graph;
    \item saturation depth as a structural quantity;
    \item reverse-degree popularity as local evidence;
    \item optional visibility constraints for controlled propagation;
    \item multiscale $K$-layer analysis of cluster persistence and merging.
\end{enumerate}

The rest of the paper is organized as follows. Section~\ref{sec:related} discusses related work. Section~\ref{sec:relational} reconstructs the original relational model. Section~\ref{sec:knn} instantiates it as a directed kNN graph. Section~\ref{sec:closure} defines the round-trip closure. Section~\ref{sec:leadership} discusses popularity and leadership. Section~\ref{sec:layers} introduces the multiscale $K$-layer representation. Section~\ref{sec:scale} reinterprets the old transmissivity--confidence argument as a scale-selection heuristic. Section~\ref{sec:protocol} proposes an experimental protocol. Section~\ref{sec:discussion} discusses limitations and open points.

\section{Related work}
\label{sec:related}

The proposed method is related to several families of clustering algorithms, but does not coincide exactly with any single one of them.

\subsection{Shared-neighbor and Jarvis--Patrick clustering}

Jarvis--Patrick clustering defines similarity using shared nearest neighbors \citep{JarvisPatrick1973}. Shared nearest neighbor clustering later became a robust strategy for detecting clusters with different shapes, sizes, and densities \citep{ErtozSteinbachKumar2003}. The present method shares with these approaches the intuition that two points can be similar because their neighborhoods overlap. However, it uses shared-neighbor information as part of a recursive closure process rather than as a static similarity score alone.

\subsection{Mutual and reciprocal nearest-neighbor methods}

Mutual kNN graphs and reciprocal nearest-neighbor relations have been used in clustering, outlier detection, and hierarchical procedures \citep{BritoChavezQuirozYukich1997,GowdaKrishna1978,SardanaBhatnagar2014,XieLeeWangZhou2023AdaptiveMKNN}. In the proposed framework, mutuality can be regarded as one possible visibility constraint, but not as the whole clustering rule. The round-trip closure can use directed edges, reverse edges, or admissible combinations of both.

\subsection{Reverse nearest neighbors and density-based methods}

Reverse nearest-neighbor methods use the number of incoming nearest-neighbor links as a density or popularity signal \citep{ChowdhuryAmorim2018}. Density-based clustering methods such as DBSCAN, OPTICS, and HDBSCAN detect clusters through density connectivity and hierarchical density estimates \citep{EsterKriegelSanderXu1996,AnkerstBreunigKriegelSander1999,CampelloMoulaviSander2013,CampelloMoulaviZimekSander2015}. The proposed method also uses reverse degree as evidence of local representativeness, but the main mechanism is graph closure and saturation rather than density thresholding.

\subsection{Graph-based and spectral clustering}

Spectral clustering builds a graph representation and uses eigenvectors of graph Laplacians to reveal cluster structure \citep{NgJordanWeiss2002,VonLuxburg2007}. The proposed method is also graph-based, but it is combinatorial and closure-driven. It does not require eigendecomposition and naturally yields a multiscale sequence of partitions as $K$ varies.

\section{From redundancy relations to round-trip closure}
\label{sec:relational}

The original 2013 model considered two finite sets,
\begin{equation}
A=\{a_1,\ldots,a_N\},\qquad B=\{b_1,\ldots,b_M\},
\end{equation}
connected by an asymmetric relation
\begin{equation}
T\subseteq A\times B.
\end{equation}
In the original terminology, members of $A$ acted as informants and members of $B$ as censors or receivers. In the present paper, this terminology is replaced by the more neutral language of directed relations and data points.

The central construction is the reflective composition
\begin{equation}
R_T = T^{-1}\circ T \subseteq A\times A.
\end{equation}
For $a\in A$, the induced reflective neighborhood is
\begin{equation}
R_T(a)=T^{-1}(T(a)).
\label{eq:basic_reflection}
\end{equation}
Thus, one first moves from $a$ to the objects in $B$ related to it, and then returns to all objects in $A$ that share those objects in $B$. This is the primitive round-trip operation.

For a subset $S\subseteq A$, the operator extends by union:
\begin{equation}
R_T(S)=\bigcup_{a\in S}R_T(a).
\end{equation}
Recursive reflection is defined by
\begin{equation}
R_T^0(S)=S,
\end{equation}
\begin{equation}
R_T^{t+1}(S)=R_T(R_T^t(S)),\qquad t\geq 0.
\label{eq:rt_iteration_general}
\end{equation}
A set saturates when
\begin{equation}
R_T^{s+1}(S)=R_T^s(S)
\end{equation}
for the first time. The integer $s$ is the saturation depth of $S$ under $R_T$, and the limit set is denoted
\begin{equation}
R_T^\infty(S)=R_T^s(S).
\end{equation}

Two objects $a,b\in A$ are disconnected under reflection if
\begin{equation}
R_T(a)\cap R_T(b)=\emptyset.
\end{equation}
They are reflectively similar if
\begin{equation}
R_T(a)\cap R_T(b)\neq\emptyset.
\end{equation}
The shared region
\begin{equation}
W(a,b)=R_T(a)\cap R_T(b)
\end{equation}
can be interpreted as an evidence zone or witness set. This is closely related to shared-neighbor reasoning, but it is defined for a general directed relation before any metric or kNN graph is introduced.

\begin{figure}[t]
\centering
\maybegraphics[width=0.7\linewidth]{fig_2013_redundancy_relation.pdf}{Schematic of the original redundancy relation $T\subseteq A\times B$ and its reflective composition $T^{-1}\circ T$.}
\caption{Relational origin of the round-trip closure. A directed relation maps objects from $A$ to $B$; the inverse relation brings the process back to $A$, producing reflective neighborhoods.}
\label{fig:relational_origin}
\end{figure}

\section{Directed kNN graph formulation}
\label{sec:knn}

Let
\begin{equation}
X=\{x_1,\ldots,x_N\}\subset\R^d
\end{equation}
be a finite dataset equipped with a metric $\delta:X\times X\to\R_{\geq 0}$. For each point $x_i$, let
\begin{equation}
\calN_K^+(x_i)=\KNN_K(x_i)
\end{equation}
denote the set of its $K$ nearest proper neighbors. The directed kNN relation is
\begin{equation}
T_K\subseteq X\times X,
\end{equation}
where
\begin{equation}
(x_i,x_j)\in T_K \quad\Longleftrightarrow\quad x_j\in\calN_K^+(x_i).
\end{equation}
The corresponding directed graph is
\begin{equation}
\calG_K=(X,E_K),
\end{equation}
with
\begin{equation}
E_K=\{(x_i,x_j):x_j\in\calN_K^+(x_i)\}.
\end{equation}

The reverse neighborhood of $x_j$ is
\begin{equation}
\calN_K^-(x_j)=T_K^{-1}(x_j)=\{x_i\in X:x_j\in\calN_K^+(x_i)\}.
\end{equation}
The reverse degree, or popularity, is
\begin{equation}
\rdeg_K(x_j)=|\calN_K^-(x_j)|.
\label{eq:reverse_degree}
\end{equation}
All vertices have out-degree $K$ in the ideal case, but the reverse degree varies and encodes how often a point is selected by other points. In the 2018 formulation, this quantity was called receptivity.

\begin{figure}[t]
\centering
\maybegraphics[width=0.72\linewidth]{fig_2018_knn_graph.pdf}{Directed 2-nearest-neighbor graph from the earlier kNN redundancy-network manuscript.}
\caption{Directed kNN graph. Each point emits $K$ outgoing links, whereas the number of incoming links varies. This asymmetry is the source of the popularity signal used by the method.}
\label{fig:knn_graph}
\end{figure}

\section{Round-trip closure in a kNN graph}
\label{sec:closure}

The kNN instance of the reflective operator is
\begin{equation}
R_K = T_K^{-1}\circ T_K.
\end{equation}
For a point $x_i$, this gives
\begin{equation}
R_K(x_i)=T_K^{-1}(T_K(x_i)).
\label{eq:knn_reflection}
\end{equation}
Equivalently,
\begin{equation}
R_K(x_i)=\bigcup_{x_j\in\calN_K^+(x_i)}\calN_K^-(x_j).
\end{equation}
Thus, $R_K(x_i)$ is the set of points that share at least one of $x_i$'s outgoing kNN targets.

For a subset $S\subseteq X$,
\begin{equation}
R_K(S)=\bigcup_{x_i\in S}R_K(x_i).
\end{equation}
The recursive round-trip expansion is
\begin{equation}
R_K^0(S)=S,
\end{equation}
\begin{equation}
R_K^{t+1}(S)=R_K(R_K^t(S)).
\label{eq:knn_rt_iteration}
\end{equation}
The saturated closure is
\begin{equation}
R_K^\infty(S)=R_K^s(S),\qquad R_K^{s+1}(S)=R_K^s(S).
\end{equation}

A more permissive forward-expansion variant can also be defined by
\begin{equation}
F_K(S)=S\cup\bigcup_{x_i\in S}\calN_K^+(x_i),
\end{equation}
and iterated until saturation. In implementation, the distinction between forward expansion, reflective expansion, and constrained expansion must be specified explicitly. The present paper uses \emph{round-trip closure} to refer to the reflective family of operators derived from $T_K^{-1}\circ T_K$ and its constrained variants.

\subsection{Saturated closures as cluster evidence}

The saturated closure $R_K^\infty(x_i)$ is not merely a local neighborhood. It is the stable set reached by repeated reflective propagation from $x_i$. If two points generate identical or strongly overlapping saturated closures, they provide evidence for belonging to the same cluster.

A strict equivalence relation can be defined by
\begin{equation}
x_i\sim_K x_j
\quad\Longleftrightarrow\quad
R_K^\infty(x_i)=R_K^\infty(x_j).
\end{equation}
A relaxed version uses the Jaccard overlap
\begin{equation}
J_K(x_i,x_j)=
\frac{|R_K^\infty(x_i)\cap R_K^\infty(x_j)|}
{|R_K^\infty(x_i)\cup R_K^\infty(x_j)|},
\end{equation}
and declares equivalence when
\begin{equation}
J_K(x_i,x_j)\geq \theta.
\end{equation}
The threshold-free implementation may instead identify absorbing leaders or connected components in the induced closure graph.

\section{Popularity, leadership, and visibility}
\label{sec:leadership}

The 2018 kNN redundancy-network formulation introduced a leadership mechanism based on reverse degree. Given the outgoing neighbors of $x_i$, a leader candidate is selected as the most popular neighbor:
\begin{equation}
\phi_K(x_i)=\argmax_{x_j\in\calN_K^+(x_i)}\rdeg_K(x_j).
\label{eq:leader_function}
\end{equation}
The image
\begin{equation}
E_K=\phi_K(X)
\end{equation}
is the leadership constellation. Points in $E_K$ act as locally preferred representatives because many points select them as neighbors.

The leadership function defines a directed graph in which every point points to one selected leader. This graph can generate a forest of attraction toward local representatives. In earlier experiments, this mechanism was used to parse hierarchical clusters in small kNN examples.

\begin{figure}[t]
\centering
\maybegraphics[width=0.72\linewidth]{fig_2018_leadership_graph.pdf}{Leadership graph or forest extracted from the earlier kNN redundancy-network manuscript.}
\caption{Leadership by popularity. Each point selects the most popular among its $K$ nearest neighbors. Cluster roots emerge as local representatives with high reverse degree.}
\label{fig:leadership_graph}
\end{figure}

\subsection{Visibility constraints}

Popularity alone can over-connect regions when $K$ is too large or when bridges exist between clusters. Therefore, the modern formulation allows a visibility function
\begin{equation}
\nu_K:X\times X\to\{0,1\}
\end{equation}
that controls which directed links are admissible during expansion. A constrained neighborhood is
\begin{equation}
\calN_{K,\nu}^+(x_i)=\{x_j\in\calN_K^+(x_i):\nu_K(x_i,x_j)=1\}.
\end{equation}
Possible visibility criteria include:
\begin{itemize}
    \item mutuality: $x_i\in\calN_K^+(x_j)$;
    \item popularity dominance: $\rdeg_K(x_j)\geq \rdeg_K(x_i)$;
    \item local distance normalization;
    \item anisotropic Mahalanobis neighborhoods;
    \item stability under successive $K$ layers.
\end{itemize}

The historical term \emph{Tallyho} may be used as the name of a particular visibility/popularity rule, but the article should define the rule mathematically before relying on the name.

\section{Cluster partitions and multiscale $K$-layers}
\label{sec:layers}

For each value of $K$, the algorithm produces a partition or cover
\begin{equation}
\calP_K=\{C_1^{(K)},\ldots,C_{m_K}^{(K)}\}.
\end{equation}
The number of clusters $m_K$ is not fixed in advance. It emerges from the closure and merging rules. As $K$ increases, small components are expected to merge, although strict monotonicity is not guaranteed under all visibility rules.

To compare two consecutive layers $K$ and $K+1$, define the transition matrix
\begin{equation}
W_{ab}^{(K,K+1)}=|C_a^{(K)}\cap C_b^{(K+1)}|.
\label{eq:transition_weight}
\end{equation}
A row-normalized transition is
\begin{equation}
P_{ab}^{(K,K+1)}=
\frac{|C_a^{(K)}\cap C_b^{(K+1)}|}{|C_a^{(K)}|},
\end{equation}
and a symmetric overlap score is
\begin{equation}
J_{ab}^{(K,K+1)}=
\frac{|C_a^{(K)}\cap C_b^{(K+1)}|}{|C_a^{(K)}\cup C_b^{(K+1)}|}.
\end{equation}
These quantities define a cluster transition diagram across $K$ layers. Persistent clusters appear as stable paths across layers; unstable fragments disappear quickly; mergers appear as many-to-one transitions.

\begin{figure}[t]
\centering
\maybegraphics[width=0.78\linewidth]{fig_k_layer_transition.pdf}{Diagram of cluster transitions across increasing $K$ layers. Node sizes represent cluster populations; edge widths represent shared points between consecutive layers.}
\caption{Multiscale $K$-layer analysis. Each layer is a partition obtained for one value of $K$. Edges encode population overlap between clusters in consecutive layers.}
\label{fig:k_layers}
\end{figure}

\section{Transmissivity, confidence, and scale selection}
\label{sec:scale}

The 2013 formulation introduced two competing quantities: transmissivity and confidence. In the original notation, if $m$ is the average number of outgoing connections from $A$ to $B$ and $n$ is the average number of reverse connections from $B$ to $A$, the expected transmissivity was expressed as
\begin{equation}
\tau=\frac{mn}{MN}.
\end{equation}
A complementary confidence quantity was approximated by
\begin{equation}
\gamma\approx\frac{1}{mn}.
\end{equation}
Equating transmissivity and confidence yields
\begin{equation}
(mn)^2=MN.
\end{equation}

In the present context, this argument can be reinterpreted as a scale-selection heuristic. Increasing $K$ increases propagation capacity, but also dilutes locality. Small $K$ values preserve local selectivity but may fragment clusters. Large $K$ values improve connectivity but may over-saturate the graph and merge distinct structures. Informative cluster scales are expected to occur in intermediate regimes where propagation and selectivity are balanced.

For a directed kNN graph, the total number of outgoing links is approximately $NK$, so the normalized graph density is
\begin{equation}
\rho_K=\frac{|E_K|}{N(N-1)}\approx\frac{K}{N-1}.
\end{equation}
Since the mean reverse degree is also approximately $K$, the variance and concentration of reverse degree may be more informative than its mean:
\begin{equation}
\operatorname{Var}(\rdeg_K),\qquad
\max_i\rdeg_K(x_i),\qquad
H(\rdeg_K).
\end{equation}
These quantities can be used together with saturation depth and partition stability to identify meaningful $K$ ranges.

\section{Algorithmic description}
\label{sec:algorithm}

The following pseudocode summarizes the method at a conceptual level.

\begin{verbatim}
Input:
    X          dataset with N data points
    K_values   neighborhood sizes to be explored
    visibility optional admissibility rule
Output:
    P_K        cluster partition for each K

for K in K_values:
    build directed kNN graph G_K
    compute reverse degrees rdeg_K

    for each point x_i in X:
        S = {x_i}
        repeat:
            S_old = S
            targets = outgoing neighbors of points in S
            if visibility is enabled:
                keep only visible/admissible targets
            S = inverse-neighborhood of targets
            S = S union S_old
        until S == S_old
        store saturated closure R_K^infty(x_i) = S

    merge points according to closure equivalence,
    closure overlap, or absorbing leader rules

    output partition P_K

compute transition matrices between consecutive K layers
\end{verbatim}

This pseudocode intentionally separates the core closure mechanism from optional rules. A reproducible implementation must specify: the kNN construction method, tie handling, visibility rule, leader election rule, closure-equivalence criterion, and merge policy.

\section{Experimental protocol}
\label{sec:protocol}

The method should be evaluated using both synthetic and real datasets. The first paper should prioritize controlled synthetic datasets, because they reveal the geometric behavior of the algorithm.

\subsection{Synthetic datasets}

Recommended synthetic benchmarks include:
\begin{itemize}
    \item Gaussian blobs with equal density;
    \item blobs with variable density;
    \item two moons;
    \item concentric circles;
    \item anisotropic Gaussian clusters;
    \item Swiss roll or manifold-type data;
    \item uniform/Poisson point clouds;
    \item clusters connected by sparse bridges.
\end{itemize}

Each dataset tests a different aspect of the method. Variable-density blobs test scale adaptation. Moons and circles test nonconvex structure. Uniform data test expected saturation in homogeneous media. Bridge datasets test resistance to spurious merging.

\subsection{Baselines}

The initial comparison should include at least:
\begin{itemize}
    \item k-means \citep{MacQueen1967,Lloyd1982};
    \item DBSCAN \citep{EsterKriegelSanderXu1996};
    \item HDBSCAN \citep{CampelloMoulaviSander2013,CampelloMoulaviZimekSander2015};
    \item spectral clustering \citep{NgJordanWeiss2002,VonLuxburg2007};
    \item shared-nearest-neighbor clustering \citep{JarvisPatrick1973,ErtozSteinbachKumar2003};
    \item mutual kNN graph clustering when available.
\end{itemize}

\subsection{External and internal metrics}

When ground-truth labels are available, use adjusted Rand index and normalized mutual information \citep{HubertArabie1985,VinhEppsBailey2010}. Cluster IDs must not be compared directly with class labels, because cluster labels are arbitrary up to permutation.

Internal metrics may include silhouette, Davies--Bouldin, and Calinski--Harabasz indices \citep{Rousseeuw1987,DaviesBouldin1979,CalinskiHarabasz1974}. These metrics should be interpreted cautiously for nonconvex structures, where Euclidean compactness may penalize topologically correct clusters.

\subsection{Method-specific diagnostics}

The method should also report diagnostics tied to its own mechanism:
\begin{itemize}
    \item number of clusters as a function of $K$;
    \item cluster-size distribution by $K$;
    \item average and maximum saturation depth;
    \item reverse-degree distribution;
    \item persistence of clusters across $K$ layers;
    \item transition matrices $W^{(K,K+1)}$;
    \item fraction of points changing cluster between layers;
    \item sensitivity to visibility rules.
\end{itemize}

\section{Expected behavior}

The expected behavior can be summarized as follows.

For small $K$, the graph may be fragmented and closure regions may be small. This regime is useful for identifying local microstructure but may oversegment the data. For intermediate $K$, clusters corresponding to meaningful data structures should persist across several layers. For large $K$, round-trip propagation may saturate large connected regions, possibly merging distinct clusters. In homogeneous point clouds, full saturation into a single cluster is not necessarily a failure; it may reflect the absence of persistent internal structure at the inspected scale.

This interpretation is particularly important for scientific datasets where clusters may be scale-dependent rather than absolute. The algorithm should therefore be evaluated not only by a single best partition, but also by the stability and interpretability of the full $K$-layer sequence.

\section{Discussion and open issues}
\label{sec:discussion}

The present formulation still leaves several technical choices to be finalized.

First, the exact closure operator must be fixed. The pure relational operator $T_K^{-1}\circ T_K$ differs from forward neighbor expansion. Both are meaningful, but they yield different dynamics. A final manuscript should state which one is implemented and why.

Second, the role of popularity and leadership must be separated from the role of closure. Popularity can define leaders, guide visibility, or break ties, but the core clustering mechanism should not be confused with simple weighted voting.

Third, visibility constraints need formal definition. Without such constraints, increasing $K$ can produce rapid saturation. With overly restrictive constraints, the method may fragment. The correct balance is part of the model design.

Fourth, computational complexity must be analyzed. Building the kNN graph can dominate runtime. Once the graph is built, closure operations can be optimized using breadth-first search, memoization of saturated closures, component compression, and union--find structures.

Finally, the relation to existing methods must be stated conservatively. The method should not be presented as the first nearest-neighbor clustering algorithm. Its defensible novelty lies in recursive round-trip closure, saturation analysis, popularity/visibility constraints, and multiscale $K$-layer interpretation.

\section{Conclusion}

This paper formulated a clustering method based on round-trip closure in directed kNN graphs. The method can be understood as a modern instantiation of an earlier relational redundancy model, where the reflective composition $T^{-1}\circ T$ defined redundancy neighborhoods and recursive saturation. In the kNN setting, this becomes a graph-topological mechanism for expanding local evidence until stable regions emerge.

The proposed framework combines directed kNN topology, reverse-degree popularity, optional visibility constraints, recursive closure, saturation depth, and multiscale analysis across $K$ layers. Its strongest use case is not merely producing a single clustering output, but revealing how cluster structure forms, persists, and merges as the local neighborhood scale changes.

Future work should complete the formal implementation specification, validate the method on controlled synthetic datasets, compare it with established baselines, and apply it to scientific data where multiscale graph topology may expose structure not captured by centroid-based or density-threshold methods alone.

\bibliographystyle{plainnat}
\bibliography{round_trip_knn_refs}

\end{document}
